
Une étudiante amateure de « bungee » oscille à la verticale au-dessus d'une rivière. Sa position en fonction du temps est donnée par :
\[ x(t) = \num{\v{Amp}} \sin\left(\num{\v{w}}\, t\right) \]
où $x$ est en mètres et $t$ en secondes.

La masse de l'étudiante est de \qty{\v{M}}{kg}. Lorsque l'étudiante est à la position la plus basse de son oscillation, elle se trouve à \qty{\v{H}}{m} au-dessus de la surface de l'eau.

\begin{enumerate}[label=\alph*)]
    \item On suppose que l'élastique du bungee se comporte comme un ressort. Quelle est sa constante de rappel ?
    \item Pendant chaque cycle de l'oscillation, durant combien de temps l'étudiante se trouve-t-elle à moins de deux mètres de la surface de l'eau ? Tracez un graphique approximatif de la position en fonction du temps et identifiez sur votre graphique les points utilisés pour le calcul.
    \item Combien de fois par cycle l'énergie cinétique de l'étudiante est-elle maximale ? À quel(s) moment(s) ? Expliquez votre réponse qualitativement, sans aucun calcul.
\end{enumerate}